% ===== 第二篇扉页 + 第 4 章 §4.1–§4.3（印刷 pp.131–146 / PDF p.146–161）· tr-01 =====
% 块首按 BRIEF 写 \part{第二篇\quad 例子：$\PP^2$}（原文 p.131 篇扉页）与 \section{Mikhalkin 曲线计数公式}（第 4 章标题）。
% PDF p.147（印刷 p.132）为空白页，不产生内容。
% tag 区间：4.1（p.136 式 ∂X_0 := ⋃D_ω）、4.2（p.139 交换图）、4.3（p.140 ghost sheaf 茎上的交换图）、
%            4.4（p.140 两行 f^*(z^{α_i})）、4.5（p.140 三行 f̄^#）。
% 本块止于印刷 p.146 末“……pre-log curves:”；p.147 起 Definition 4.21 属 c04b（tr-03）。
% 交换图：(4.2)(4.3) 及两处未编号小交换图按 lead 裁决用 array 排；p.135 Figure 1 与 p.137 的“证明结构图”、
% p.146 的大交换图按原图裁剪（f4_1.png / f4_diagram.png / f4_diagram2.png）。

\part{第二篇\quad 例子：$\PP^2$}
% [原文 p.131] 第二篇扉页：Example: $\PP^2$（PDF p.146）。
% [原文 p.132] 空白页（PDF p.147）。

% [原文 p.133]
\section{Mikhalkin 曲线计数公式}

\subsection{定理的陈述与证明纲要}
我们将先给出 Mikhalkin 定理的陈述，其证明将在本章余下部分给出。

照例固定 $M = \Z^n$、$N = \operatorname{Hom}_\Z(M,\Z)$ 等等，并固定 $M_\R$ 中的一个完整扇 $\Sigma$。沿用 §1.3
的记号，这给出格 $T_\Sigma$ 与映射 $r : T_\Sigma \to M$；在 $r$ 的核中取
\[
\Delta = \sum_{\rho\in\Sigma^{[1]}} d_\rho t_\rho.
\]
若 $n = 2$，则 §1.3 中已定义了数 $N^{0,\mathrm{trop}}_{\Delta,\Sigma}$，它（带重数地）计数模空间
$\cM_{0,|\Delta|-1}(\Sigma,\Delta)$ 中经过 $|\Delta|-1$ 个一般点的热带曲线。

我们首先需要定义与之对应的全纯的或代数几何的计数。我们将如第 2 章那样计数稳定曲线，但对所考虑的曲线加一些
限制。此外还需要指定曲线所在的同调类。对第一点，我们用

\begin{definition}
若曲线 $C \subseteq X_\Sigma$ 与环面簇 $X_\Sigma$ 的所有余维数 $>1$ 的环面层都不相交，则称它是\emph{环面横截的}。
（特别地，$C$ 不含任何位于余维数 $1$ 的层中的不可约分支，因为那样的分支还会与某个更高余维数的层相交。）

若 $f : C \to X_\Sigma$ 是稳定映射，则称它是\emph{环面横截的}，如果 $f(C) \subseteq X_\Sigma$ 环面横截，且 $C$ 的
任何不可约分支都不映入 $\partial X_\Sigma$。\hfill$\square$
\end{definition}

% [Tutor Note: 原文 p.133，定义 4.1 末尾的 $\square$ 是原文本身的排版痕迹（定义后误留了证明结束符），照录不改。
%  “环面横截”的关键在于：$C$ 不碰到余维 $\ge2$ 的环面层，因此 $C$ 与每个除子 $D_\rho$ 的交都是横截的、且落在
%  该除子的光滑部分；这一段是后面 Proposition 4.2、4.9 全部论证的前提。术语按 glossary 定为“环面横截”。]

对次数，我们用如下结果：

\begin{proposition}
若 $X_\Sigma$ 非奇异，则
\[
H_2(X_\Sigma,\Z) \cong \ker r.
\]
该同构把 $\beta \in H_2(X_\Sigma,\Z)$ 送到
\[
\sum_{\rho\in\Sigma^{[1]}} (\beta.D_\rho)t_\rho \in T_\Sigma.
\]
\end{proposition}

\begin{proof}
由式 (3.2) 可见 $\operatorname{Cl}(X_\Sigma) = \operatorname{Hom}_\Z(\ker r,\Z)$；但 $X_\Sigma$ 非奇异（且所有
环面簇都是有理的），故 $\operatorname{Cl}(X_\Sigma) \cong \operatorname{Pic}(X_\Sigma) \cong H^2(X_\Sigma,\Z)$。
于是 $\ker r \cong H_2(X_\Sigma,\Z)$。更具体地说，$T_\Sigma^\vee$ 中形如 $\sum a_\rho t_\rho^*$ 的元素代表
$\operatorname{Cl}(X_\Sigma)$ 中由 $\sum a_\rho D_\rho$ 给出的除子。此外，使 $H_2(X_\Sigma,\Z)$ 成为
$H^2(X_\Sigma,\Z)$ 的对偶的配对 $H_2(X_\Sigma,\Z)\times H^2(X_\Sigma,\Z) \to \Z$ 由相交配对
$(\beta,D_\rho) \mapsto \beta\cdot D_\rho$ 给出。这就给出该同构的显式描述。
\end{proof}

于是 $\Delta \in \ker r$ 确定了 $\Delta \in H_2(X_\Sigma,\Z)$。因此如第 2 章那样，我们可以用
$\overline{\cM}_{0,|\Delta|-1}(X_\Sigma,\Delta)$ 记亏格零、带 $|\Delta|-1$ 个标记点、代表同调类 $\Delta$ 的
稳定映射映入 $X_\Sigma$ 的模空间。

% [原文 p.134]
\begin{definition}
设 $Q_1,\dots,Q_{|\Delta|-1} \in X_\Sigma$ 是一般点。定义
\[
N^{0,\mathrm{hol}}_{\Delta,\Sigma}
= \#\left\{f\in \overline{\cM}_{0,|\Delta|-1}(X_\Sigma,\Delta)
\ \middle|\
\begin{array}{l}
f : (C,x_1,\dots,x_{|\Delta|-1}) \to X_\Sigma \text{ 环面横截，}\\
\text{且对 } 1\le i\le|\Delta|-1 \text{ 有 } f(x_i) = Q_i
\end{array}\right\}.
\]
\end{definition}

目前还不清楚这是一个有限数，但这将在证明主定理的过程中得到证明。不过一旦知道了这一点，也就知道它不依赖于
$Q_1,\dots,Q_{|\Delta|-1}$ 的选取（只要它们取得一般），因为随着 $Q_1,\dots,Q_{|\Delta|-1}$ 变动，被计数的集合
构成一个代数族，所以对任意两个一般选取，其基数相同。于是我们可以陈述主定理：

\begin{theorem}
设 $n = 2$。则 $N^{0,\mathrm{hol}}_{\Delta,\Sigma}$ 有限，且
\[
N^{0,\mathrm{trop}}_{\Delta,\Sigma} = N^{0,\mathrm{hol}}_{\Delta,\Sigma}.
\]
\end{theorem}

注意我们还不知道 $N^{0,\mathrm{trop}}_{\Delta,\Sigma}$ 不依赖于 $M_\R$ 中点的选取；这一点由上一定理推出。不过
下一章我们将给出这一独立性的另一个证明，它完全是组合的。

我们给出的证明不是 Mikhalkin 原来的证明，而是 Siebert 与 Nishinou [86] 的证明。后者的证明在所有维数都适用，
而 Mikhalkin 原来的证明只在二维适用。不过我们将限制在二维。高维情形中 $N^{0,\mathrm{trop}}_{\Delta,\Sigma}$ 的
定义稍为复杂一些，证明中涉及的组合也复杂一些。但限制到二维其实并不能省下多少，所以读完本章之后，查阅 [86] 的
一般情形应当很容易。

基本想法是构造 $X_\Sigma$ 的一个单参数退化，其风味与例 3.6 非常相似。这个退化将针对用来定义
$N^{0,\mathrm{trop}}_{\Delta,\Sigma}$ 的那组特定点 $P_1,\dots,P_{|\Delta|-1} \in M_\R$ 作适配。然后我们将在
以下几者之间建立对应关系：$M_\R$ 中经过 $P_1,\dots,P_{|\Delta|-1}$ 的热带曲线、该退化中心纤维中的“对数稳定
曲线”、以及该退化一般纤维中的普通稳定曲线。

更具体地说，取点 $P_1,\dots,P_s \in M_{\Q} = M\otimes_\Z\Q$，其中 $s = |\Delta|-1$。注意这里我们只考虑坐标
为有理数（而非实数）的点。这并不改变前面的讨论：引理 1.20 的一般位置论证在 $\Q$ 上同样成立，所以可以假定这些
点的选取使得 $X_\Sigma$ 中满足 $h(x_i)=P_i$ 的带标记亏格零热带曲线
$h : (\Gamma,x_1,\dots,x_{|\Delta|-1}) \to M_\R$ 只有有限多条，且这些曲线都是单的。用
\[
\cM_{0,s}(\Sigma,\Delta,P_1,\dots,P_s)
\]
记这些曲线之集。

\begin{definition}
给定 $\Sigma$ 与一般点 $P_1,\dots,P_{|\Delta|-1} \in M_{\Q}$（$\dim M_{\Q} = 2$），称 $M_\R$ 的一个有限多面体
分解 $\cP$ 是\emph{好的}，如果它满足下列性质：
\begin{enumerate}[label=(\arabic*),leftmargin=*]
\item 对 $\sigma\in\cP$，$\sigma$ 的面具有有理斜率，且顶点都在 $M_{\Q}$ 中。此外每个 $\sigma\in\cP$ 至少
有一个顶点。
\item 若 $\sigma\in\cP$，则 $\sigma$ 的渐近锥 $\operatorname{Asym}(\sigma)$ 是扇 $\Sigma$ 的元素。此外 $\Sigma$ 的
每个锥都作为某个 $\sigma\in\cP$ 的渐近锥出现。
\item 任何 $h \in \cM_{0,s}(\Sigma,\Delta,P_1,\dots,P_s)$ 的像都落在 $\cP$ 的一维骨架中。
\item 每个 $P_i$ 都是 $\cP$ 的顶点。
\end{enumerate}
\end{definition}

% [原文 p.135]
\begin{figure}[htbp]
\centering
\includegraphics[width=0.92\textwidth]{f4_1.png}
\caption*{\small\textbf{图1}}
\end{figure}

% [Tutor Note: 原文 p.135，图 1。左图是 $\PP^2$ 中一条次数为 $2$ 的热带曲线（含 5 个标记点），右图是命题 4.6
%  构造出的多面体分解 $\cP$ 的样子：注意 $\cP$ 的一维骨架包含整条热带曲线，而每个 $P_i$ 都是 $\cP$ 的顶点。
%  原书此图没有说明文字，只有编号“Figure 1”，故这里只复现编号。图在原文中排在第 (2)、(3) 条之间（页顶浮动），
%  因 LaTeX 浮动体不能放进 definition 环境内部，这里把定义 (1)–(4) 排完再插图，语义不受影响。]

\begin{proposition}
好的多面体分解存在。
\end{proposition}

\begin{proof}
若 $\cP$ 与 $\cP'$ 是两个满足条件 (1) 与 (2) 的多面体分解，则
\[
\cP'' = \{\sigma\cap\sigma' \mid \sigma\in\cP,\ \sigma'\in\cP'\}
\]
也满足这两个条件。我们如下使用这一点。取 $\cP$ 为这样的多面体分解：其极大胞腔是 $M_\R\setminus S$ 的连通
分支的闭包，其中 $S$ 是 $\cM_{0,s}(\Sigma,\Delta,P_1,\dots,P_s)$ 中参数化热带曲线的像与 $\Sigma^{[1]}$ 的元素
之并。（若对一切 $\rho$ 都有 $d_\rho>0$，其中 $\Delta=\sum_\rho d_\rho t_\rho$，则后一部分并不需要。）这个
$\cP$ 满足 (1)、(2)、(3)。取 $\cP'$ 为这样的多面体分解：其极大胞腔是 $M_\R\setminus S'$ 的连通分支的闭包，
其中
\[
S' = \bigcup_{\rho\in\Sigma^{[1]}} \bigcup_{i=1}^s (P_i+\rho).
\]
则 $\cP'$ 满足条件 (1)、(2)、(4)。上面给出的 $\cP$ 与 $\cP'$ 的加细 $\cP''$ 满足条件 (1)–(4)。
\end{proof}

例如，若我们在 $M_\R$ 中取落在图 1 左侧所画的那条 $\PP^2$ 中次数为 $2$ 的热带曲线上的五个点，则可以按右侧
所画的那样取 $\cP$。

不失一般性，我们其实可以假定 $\cP$ 是格多面体分解：把 $M$ 换成超格 $\frac{1}{p}M$，其中 $p$ 是 $\cP$ 所有顶点
坐标的公分母。这当然使点 $P_1,\dots,P_s$ 也落在 $M$ 中，但不改变数 $N^{0,\mathrm{trop}}_{\Delta,\Sigma}$。在
余下的讨论中我们固定这个好的 $\cP$。

现在我们引入记号
\[
\widetilde{M} = M\oplus\Z,\qquad \widetilde{N} = \operatorname{Hom}_\Z(\widetilde{M},\Z),
\]
并如下定义 $\widetilde{M}_\R$ 中的扇 $\Sigma_{\cP}$。对 $\sigma\in\cP$，令
\[
C(\sigma) = \{(rm,r) \mid r\ge0,\ m\in\sigma\} \subseteq \widetilde{M}_\R
\]

% [原文 p.136]
为 $\sigma$ 上的锥，如 §3.1.3 所述。然后取 $\Sigma_{\cP}$ 为所有 $C(\sigma)$（$\sigma\in\cP$）的面构成的扇。
容易看出 $\Sigma_{\cP}$ 是扇。此外注意 $C(\sigma)\cap(M_\R\oplus\{0\}) = \operatorname{Asym}(\sigma)$，如
§3.1.3 所述。又由定义 4.5 的性质 (2)，$\operatorname{Asym}(\sigma) \in \Sigma$。并且 $\Sigma$ 的每个锥都这样
出现。因此
\[
\Sigma = \{\tau\in\Sigma_{\cP} \mid \tau \subseteq M_\R\oplus\{0\}\}.
\]

把这些放在一起，我们得到一个环面簇
\[
X := X_{\Sigma_{\cP}},
\]
它带有由单项式 $z^{(0,1)}$（$(0,1)\in\widetilde{N} = N\oplus\Z$）诱导的正则函数
$\pi : X \to \mathbb{A}^1_{k}$。注意 $\pi^{-1}(0)$ 是那些使 $z^{(0,1)}$ 为零的环面除子之并。这些正是对应于
$\Sigma_{\cP}$ 中形如 $C(v)$（$v\in\cP$ 为顶点）的锥的除子。特别地 $\pi^{-1}(0)$ 是若干环面簇之并。另一方面，
若把这些除子从 $X_{\Sigma_{\cP}}$ 中去掉，所得到的环面簇对应的扇 $\Sigma$ 不是放在 $M_\R$ 中，而是放在
$\widetilde{M}_\R$ 中。于是
\[
X \setminus \pi^{-1}(0) \cong X_\Sigma \times \mathbb{G}_m.
\]

退化 $\pi : X \to \mathbb{A}^1_{k}$ 就是我们要研究的 $X_\Sigma$ 的退化。

对 $X_0 := \pi^{-1}(0)$ 的环面层，我们采用如下记号。每个胞腔 $\tau\in\cP$ 给出锥
$C(\tau)\in\Sigma_{\cP}$，从而给出环面层 $D_{C(\tau)}$。该层含于 $X_0$。为简化记号，我们把它写作 $D_\tau$。
定义 $D_\tau$ 的扇记作 $\Sigma_\tau$，它位于 $M_\R/T_\tau\tau$ 中。用 $\cP$ 中含 $\tau$ 的多面体可把它定义为
\[
\Sigma_\tau := \{T_\tau\sigma/T_\tau\tau \mid \sigma\in\cP,\ \tau\subseteq\sigma\}.
\]

用 $\partial X_0$ 记 $X_0$ 中不含于 $\operatorname{Sing}(X_0)$ 的一维环面层之并。这正好就是并
\begin{equation}
\partial X_0 := \bigcup_{\substack{\omega\in\cP^{[1]}\\ \omega\ \text{非紧}}} D_\omega, \tag{4.1}
\end{equation}
因为正是那些非紧的 $\omega$ 对应的 $D_\omega$ 不含于 $X_0$ 的两个不可约分支之中。也可以写成
\[
\partial X_0 = \overline{\partial X \setminus X_0}\cap X_0.
\]

% [Tutor Note: 原文 p.136，式 (4.1)。核对锚点：$X_0=\pi^{-1}(0)$ 的不可约分支一一对应于 $\cP$ 的顶点 $v$
%  ——即对应于 $\Sigma_{\cP}$ 中形如 $C(v)$ 的锥；因此只有“非紧”的 $\omega$ 所对应的 $D_\omega$ 不会被两个分支
%  共用，而它们恰是 $X_0$ 的一维层中不含于 $\operatorname{Sing}(X_0)$ 的那部分。这个 $\partial X_0$ 正是后面
%  定义 4.10 (3) 里 $f^{-1}(\partial X_0)$ 所对应的无界边集合。]

接下来我们解释点 $P_1,\dots,P_s$ 的作用。令
\[
L_i = \Z(P_i,1) \subseteq \widetilde{M}
\]
为由 $(P_i,1)$ 生成的一维子格。于是有环面的包含 $\mathbb{G}(L_i) \subseteq \mathbb{G}(\widetilde{M})$，它在
$X$ 上有作用。现在在 $X$ 的开环面轨道中取点 $Q_1,\dots,Q_s$，并考虑 $X$ 的闭子集
$\overline{\mathbb{G}(L_i)\cdot Q_i}$。这里我们用到 $\mathbb{G}(\widetilde{M})$（从而是 $\mathbb{G}(L_i)$）在
$X$ 上的典范作用。

\begin{proposition}
$\pi|_{\overline{\mathbb{G}(L_i)\cdot Q_i}}$ 是 $\overline{\mathbb{G}(L_i)\cdot Q_i}$ 与
$\mathbb{A}^1_{k}$ 的同构。
\end{proposition}

\begin{proof}
注意复合 $\mathbb{G}(L_i) \to \mathbb{G}(\widetilde{M}) \xrightarrow{\pi} \mathbb{G}(\Z)$ 是同构，因为它由格
同构 $L_i \to \widetilde{M} \to \Z$ 诱导。因此 $\pi|_{\mathbb{G}(L_i)\cdot Q_i}$ 是到其像上的同构。注意
$\mathbb{G}(L_i)\cdot Q_i$ 在 $\pi^{-1}(\mathbb{A}^1_{k}\setminus\{0\})$ 中是闭的。由于 $\pi$ 是逆紧的，
取 $\mathbb{G}(L_i)\cdot Q_i$ 在 $X$ 中的闭包时，恰好多出一个落在 $\pi^{-1}(0)$ 中的点，且
$\pi|_{\overline{\mathbb{G}(L_i)\cdot Q_i}}$ 是同构。
\end{proof}

% [原文 p.137]
于是在我们的设置中，$P_1,\dots,P_s \in M$ 的选取给出 $\pi$ 的截面
\[
\sigma_1,\dots,\sigma_s : \mathbb{A}^1_{k} \to X.
\]
这样我们就得到一个单参数族：$X_\Sigma$ 中 $s$ 个点的选取，退化到 $X_0$ 中点的选取。

现在我们可以精确描述这幅图景中出现的三个世界了。

\para{热带世界。}人们考虑 $M_\R$ 中经过 $P_1,\dots,P_s$ 的带标记参数化亏格零热带曲线。

\para{对数世界。}人们考虑从带标记点 $x_1,\dots,x_s$ 的亏格零对数曲线 $C^{\dagger}$ 出发的对数态射
$f : (C^{\dagger},x_1,\dots,x_s) \to X_0^{\dagger}$，满足
\[
f(x_i) = \sigma_i(0) \in X_0,
\]
并且 $X_0$ 赋以由包含 $\partial X \subseteq X$ 而来的除子对数结构所诱导的对数结构。

\para{经典世界。}数 $N^{0,\mathrm{hol}}_{\Delta,\Sigma}$ 其实在任何特征零的代数闭域上都有意义，且不依赖于域的
选取。于是我们可以考虑域 $K = \overline{k((t))}$，即 $k((t))$ 的代数闭包。包含
$k[t] \subseteq K$ 给出映射 $\operatorname{Spec}K \to \mathbb{A}^1_{k}$，而
\[
X \times_{\mathbb{A}^1_{k}} K \cong X_\Sigma \times_k K
\]
是 $\Sigma$ 在域 $K$ 上定义的环面簇。此外，截面 $\sigma_1,\dots,\sigma_s : \mathbb{A}^1_{k} \to X$ 确定点
$\sigma_1,\dots,\sigma_s : \operatorname{Spec}K \to X_\Sigma\times_k K$。用这 $s$ 个点即得
$N^{0,\mathrm{hol}}_{\Delta,\Sigma}$。对这个计数有贡献的曲线是 $\operatorname{Spec}K$ 上的环面横截曲线。
之后将会看到，这些曲线与对数世界中的曲线之间存在自然的对应。

于是定理 4.4 的证明图如下：

\begin{center}
\includegraphics[width=0.62\textwidth]{f4_diagram.png}
\end{center}

% [Tutor Note: 原文 p.137，证明结构图（原书无编号、无图注，图内文字保持原图）。读图方向：$§4.2$ 从“对数世界”
%  指向“热带世界”（对数曲线↦热带曲线），$§4.3$ 反向（热带曲线↦对数曲线，条数为 $\operatorname{Mult}(h)$），
%  $§4.4$ 从“经典世界”指向“对数世界”，$§4.5$ 从“对数世界”指向“经典世界”（对数形变理论）。$§4.6$ 把四者
%  串起来收尾——这就是本块之后 c04b 的内容。]

换言之，我们将先在 §4.2 中说明某些对数曲线如何通过“对偶相交复形”的构造给出热带曲线。在 §4.3 中我们说明如何
从热带曲线构造对数曲线，实际上证明对给定的热带曲线 $h$，联系于它的对数曲线有 $\operatorname{Mult}(h)$ 条。
特别地，这揭示了重数的来源。

在 §4.4 中我们说明 $\operatorname{Spec}K$ 上每条环面横截曲线如何给出对数曲线；在 §4.5 中通过对数形变理论
说明每条对数曲线都可形变给出 $\operatorname{Spec}K$ 上的环面横截曲线。最后在 §4.6 中我们把这一切汇总起来，
给出定理 4.4 的证明。

% [原文 p.138]
下面进入细节。

\subsection{对数世界 $\to$ 热带世界}
本节固定上一节所描述的数据：$M_\R$ 中的二维扇 $\Sigma$、次数 $\Delta\in\ker r$、$s = |\Delta|-1$、点
$P_1,\dots,P_s \in M$、一个好的格多面体分解 $\cP$，以及一般点
$Q_1,\dots,Q_s \in \mathbb{G}(\widetilde{M})$。

由这些数据我们得到一个环面簇的退化 $\pi : X \to \mathbb{A}^1_{k}$，及截面
$\sigma_1,\dots,\sigma_s : \mathbb{A}^1_{k} \to X$，其中 $\sigma_i$ 的像是
$\overline{\mathbb{G}(L_i)\cdot Q_i}$。

照前令 $X_0 = \pi^{-1}(0)$，并赋以由 $X$ 上来自包含 $\partial X \subseteq X$ 的除子对数结构所诱导的对数结构。
按例 3.24 (3) 的方式限制 $\pi : X^{\dagger} \to (\mathbb{A}^1_{k})^{\dagger}$ 得到的态射
$\pi_0 : X_0^{\dagger} \to \operatorname{Spec}k^{\dagger}$ 是对数光滑的。

\begin{definition}
$X_0^{\dagger}$ 中的\emph{环面横截对数曲线}是指一个对数态射的交换图
\[
\begin{array}{ccc}
C^{\dagger} & \xrightarrow{\ f\ } & X_0^{\dagger} \\[4pt]
{\scriptstyle g}\big\downarrow & & \big\downarrow{\scriptstyle \pi_0} \\[4pt]
\operatorname{Spec}k^{\dagger} & \xrightarrow{\ \ =\ \ } & \operatorname{Spec}k^{\dagger}
\end{array}
\]
其中 $g$ 是对数光滑且整的，$C$ 是使 $C^{\dagger}$ 精细饱和的一条曲线，$f$ 底下的概形态射是稳定映射，且对
每个顶点 $v\in\cP$，$f^{-1}(D_v)\to D_v$ 是环面横截稳定映射。
\end{definition}

% [Tutor Note: 原文 p.138，定义 4.8。“$g$ 对数光滑且整”正是 c03 中“整对数态射”的含义：对数结构由目标
%  （这里是 $\operatorname{Spec}k^{\dagger}$）拉回，因此 $C^{\dagger}$ 的对数结构完全由 $g$ 决定；
%  “$C^{\dagger}$ 精细饱和”排除掉对数结构退化得太厉害的曲线。核对锚点：图的右下角是等号，说明 $C$ 与
%  $X_0$ 的对数结构都落在同一基础 $\operatorname{Spec}k^{\dagger}$ 上。]

下面的结果给出对数曲线的关键限制，它们使我们得以建立与热带几何的联系。

\begin{proposition}
设 $f : C^{\dagger} \to X_0^{\dagger}$ 是环面横截对数曲线。若 $x\in C$ 是使
$f(x)\in\operatorname{Sing}(X_0)$ 的点，则
\begin{enumerate}[label=(\arabic*),leftmargin=*]
\item $x$ 是 $C$ 的二重点，含于 $C$ 的两个不同不可约分支 $C_1, C_2$ 中，且
$f(C_i)\subseteq D_{v_i}$（$i=1,2$），其中 $v_1,v_2\in\cP$ 是被 $\cP$ 的一条边 $\omega$ 连接的两个不同顶点。
\item 设 $w_i$ 是 $C_i$ 在 $x$ 处与 $D_\omega\subseteq D_{v_i}$ 的相交重数。这里我们指的是正则函数
$f^*(\varphi)$ 在 $C_i$ 上 $x$ 处的零点阶，其中 $\varphi$ 是 $f(x)$ 在 $D_{v_i}$ 中某邻域内定义的正则函数，
且 $\varphi=0$ 定义 $D_\omega$。则 $w_1 = w_2$。
\item 存在 $e\ge1$ 使得
\[
\overline{\cM}_{C,x} = S_e := \mathbb{N}^2\oplus_{\mathbb{N}}\mathbb{N},
\]
其中纤维和由对角映射 $\mathbb{N}\to\mathbb{N}^2$ 与乘 $e$ 的映射 $\mathbb{N}\to\mathbb{N}$ 定义。等价地，
$S_e$ 可描述为由 $\alpha_1,\alpha_2$ 与 $\rho$ 生成、并满足关系 $\alpha_1+\alpha_2 = e\rho$ 的幺半群。
（见例 3.27。）
\item 若 $\ell$ 是 $\omega$ 的仿射长度，则 $ew_i = \ell$。
\end{enumerate}
\end{proposition}

\begin{proof}
我们先描述 $X$ 中 $f(x)$ 附近的局部情形。由环面横截性条件，$f(x)$ 不可能落在零维层中。故 $f(x)$ 必落在某个
一维层中，设为 $D_\omega$，其中 $\omega\in\cP$ 是以 $v_1,v_2$ 为端点的边。不失一般性，设
$v_1 = (0,0)\in M$，$v_2 = (\ell,0)\in M$，

% [原文 p.139]
其中 $\ell$ 是 $\omega$ 的仿射长度。则锥 $C(\omega)$ 定义 $X$ 的一个仿射开子集
$X_{C(\omega)} \subseteq X$。为方便起见，我们把这个仿射开子集写作 $X_\omega$。

注意锥 $C(\omega)^\vee$ 在 $\widetilde{N}_\R = \R^3$ 中由 $(0,\pm1,0)$、$(1,0,0)$、$(-1,0,\ell)$ 生成，且
$C(\omega)^\vee\cap\widetilde{N} = \Z\oplus S_\ell$，其中 $(0,1,0)$ 生成 $\Z$，而
$\alpha_1 = (1,0,0)$、$\alpha_2 = (-1,0,\ell)$、$\rho = (0,0,1)$ 生成 $S_\ell$。注意
$X_0\cap X_\omega = D_1\cup D_2$，即对应于 $v_1$ 与 $v_2$ 的两个除子。

把 $C$ 换成 $x$ 的一个开邻域，使其在 $f$ 下的像含于 $X_\omega$ 中。我们可以假定这个开邻域足够小，使得除 $x$
以外没有点映入 $\operatorname{Sing}(X_0)$，且 $C$ 至多有两个不可约分支。这里我们用到了 $f$ 的环面横截性，
以保证 $C$ 的任何分支都不映入 $\operatorname{Sing}(X_0)$。于是现在有交换图
\[
\begin{array}{ccccc}
C^{\dagger} & \xrightarrow{\ f\ } & (X_\omega\cap X_0)^{\dagger} & \xrightarrow{\quad} & X_\omega^{\dagger} \\[4pt]
{\scriptstyle g}\big\downarrow & & \big\downarrow{\scriptstyle \pi_0} & & \big\downarrow{\scriptstyle \pi} \\[4pt]
\operatorname{Spec}k^{\dagger} & \xrightarrow{\ \ =\ \ } & \operatorname{Spec}k^{\dagger} & \xrightarrow{\quad} & (\mathbb{A}^1_{k})^{\dagger}
\end{array}
\]
把 $f$ 换成它与包含 $(X_\omega\cap X_0)^{\dagger} \to X_\omega^{\dagger}$ 的复合，于是现在有交换图
\begin{equation}
\begin{array}{ccc}
C^{\dagger} & \xrightarrow{\ f\ } & X_\omega^{\dagger} \\[4pt]
{\scriptstyle g}\big\downarrow & & \big\downarrow{\scriptstyle \pi} \\[4pt]
\operatorname{Spec}k^{\dagger} & \xrightarrow{\quad} & (\mathbb{A}^1_{k})^{\dagger}
\end{array}
\tag{4.2}
\end{equation}

首先设 $f(C)$ 含于（比如说）$D_1$ 中，且 $f(x)\in D_1\cap D_2$。注意
$z^{\alpha_1}\in\Gamma(X_\omega,\cM_{X_\omega})$，故 $f^{\#}(z^{\alpha_1})\in\Gamma(C,\cM_C)$。另一方面，
$z^{\alpha_1}$ 在 $D_2$ 上为零而在 $D_1$ 上不为零，所以
$z^{\alpha_1}\in\Gamma(X_\omega\setminus D_2,\cO^{\times}_{X_\omega})$，从而实际上
$\alpha_C(f^{\#}z^{\alpha_1}) = f^*\alpha_X(z^{\alpha_1}) = f^*(z^{\alpha_1})
\in\Gamma(C\setminus\{x\},\cO_C^{\times})$。注意到 $f^*(z^{\alpha_1})$ 在 $x$ 处不可逆，我们断定
$f^{\#}(z^{\alpha_1})$ 在 $\overline{\cM}_C$ 中的像的支撑恰好是点 $x$。

现在由例 3.26，$C^{\dagger}$ 在 $x$ 处的行为有三种可能。在该例的情形 (1)，即 $x$ 是 $C$ 的光滑点时，
$\overline{\cM}_C$ 局部是常层 $\underline{\mathbb{N}}$，因此不存在支撑仅在 $x$ 的截面。

在情形 (2)，即 $x$ 是二重点时，$\overline{\cM}_C$ 同样没有支撑仅在 $x$ 的截面。为看出这一点，注意由例
3.26 (2) 的描述，$C$ 在 $x$ 附近局部看起来像 $\operatorname{Spec}k[u,v]/(uv)$。此外例 3.27 的描述告诉
我们对数结构：它由包含 $C \subseteq V_e$ 诱导，其中
\[
V_e = \operatorname{Spec}k[u,v,t]/(uv-t^e)
\]
对某个 $e>0$，且 $C = V(t)\subseteq V_e$。这个包含在 $V_e$ 上诱导除子对数结构，再拉回到 $C$ 上的对数结构。
由于 $\overline{\cM}_{V_e}$ 的支撑在 $C$ 上，$\overline{\cM}_C$ 有支撑仅在 $x$ 的截面只有当
$\overline{\cM}_{V_e}$ 有才行。但 $\overline{\cM}_{V_e}$ 的截面以 $V_e$ 的、支撑于 $C$ 上的 Cartier 除子为
支撑，故不可能有支撑仅在 $x$ 的截面。于是第二种情形被排除。

在例 3.26 的情形 (3)，即 $x$ 是对数标记点时，$x$ 附近局部有
$\overline{\cM}_C = \underline{\mathbb{N}}\oplus\underline{\mathbb{N}}_x$，它确实有支撑仅在 $x$ 的截面。
不过我们如下排除这一情形。首先注意由例 3.26 (3)，映射 $g : C^{\dagger} \to \operatorname{Spec}k^{\dagger}$
由某个 $\rho_C \in \Gamma(C,\cM_C)$ 确定，它映到 $\bar\rho_C \in \Gamma(C,\overline{\cM}_C)$，

% [原文 p.140]
且在局部具有 $\mathbb{N}\oplus\mathbb{N}_x$ 中 $(1,0)$ 的形式。其次，由 (4.2) 的交换性，由于 $\pi$ 由
$z^{\rho}\in\Gamma(X_\omega,\cM_{X_\omega})$ 给出，我们有 $f^{\#}z^{\rho} = \rho_C$。但
$f^{\#}(z^{\ell\rho}) = f^{\#}(z^{\alpha_1}z^{\alpha_2})$。考虑 $f^{\#}(z^{\alpha_1})$ 在
$\overline{\cM}_{C,\bar x}$ 中的像 $\overline{f^{\#}(z^{\alpha_1})}$。由于 $f^{\#}(z^{\alpha_1})$ 的支撑在 $x$
且非零，必有
\[
\overline{f^{\#}(\alpha_1)} = (0,a)
\]
对某个 $a>0$。那么 $f^{\#}(z^{\alpha_2})$ 在 $\overline{\cM}_{C,\bar x}$ 中的像
$\overline{f^{\#}(z^{\alpha_2})}$ 便无从选取，使
\[
(\ell,0) = \ell\bar\rho_C = (0,a)+\overline{f^{\#}(z^{\alpha_2})}
\]
在 $\overline{\cM}_{C,\bar x}$ 中成立。这就排除了例 3.26 的第三种情形。

总之，$f(C)$ 既不含于 $D_1$ 也不含于 $D_2$。唯一可能的办法是：$x$ 是 $C$ 的二重点，是两个不可约分支
$C_1$ 与 $C_2$ 的交点，且 $f(C_1)\subseteq D_1$、$f(C_2)\subseteq D_2$。这就证明了 (1)。

对 (2)–(4)，注意 $\overline{\cM}_{C,\bar x} = S_e$ 对某个 $e\ge0$ 成立。图 (4.2) 在鬼层茎的层面是
\begin{equation}
\begin{array}{ccc}
S_e = \overline{\cM}_{C,\bar x} & \xleftarrow{\ \ \bar f^{\#}\ \ } & \overline{\cM}_{X_\omega,f(x)} = S_\ell \\[8pt]
{\scriptstyle 1\mapsto\bar\rho_C}\big\uparrow & & \big\uparrow{\scriptstyle 1\mapsto\rho} \\[8pt]
\mathbb{N} & \xleftarrow{\ \ \ =\ \ \ } & \overline{\cM}_{\mathbb{A}^1_{k},\bar 0} = \mathbb{N}
\end{array}
\tag{4.3}
\end{equation}
设 $w_1,w_2$ 是 (2) 的陈述中所定义的整数。注意在由例 3.26 (2) 给出的 $C$ 在 $x$ 附近的局部描述中，取
$C_1 = V(v)$、$C_2 = V(u)$，有
\begin{equation}
\begin{aligned}
f^*(z^{\alpha_1}) &= u^{w_1}\cdot\varphi_1\\
f^*(z^{\alpha_2}) &= v^{w_2}\cdot\varphi_2
\end{aligned}
\tag{4.4}
\end{equation}
其中 $\varphi_1,\varphi_2$ 是 $C$ 上 $x$ 附近的可逆函数。事实上，这正是 $w_1$ 与 $w_2$ 的定义。

让我们用这些信息来计算 $f^{\#}(z^{\alpha_1})$ 与 $f^{\#}(z^{\alpha_2})$。给定例 3.26 (2) 中给出的、定义 $C$ 上
对数结构的卡的描述，$f^{\#}(z^{\alpha_i})$ 可以在 $x$ 附近局部地写成 $\underline{S}_e\oplus\cO_C^{\times}$ 的
截面，比如 $(((a_i,b_i),q_i),\psi_i)$，满足
\[
f^*(z^{\alpha_i}) = \alpha_C f^{\#}(z^{\alpha_i}) =
\begin{cases}
u^{a_i}v^{b_i}\psi_i & q_i = 0\\
0 & q_i \ne 0
\end{cases}
\]
于是实际上由 (4.4) 得
\[
\bigl((a_1,b_1),q_1\bigr) = \bigl((w_1,0),0\bigr),\qquad
\bigl((a_2,b_2),q_2\bigr) = \bigl((0,w_2),0\bigr).
\]
这说明 (4.3) 中的映射 $\bar f^{\#}$ 必须满足
\begin{equation}
\begin{aligned}
\bar f^{\#}\bigl((1,0),0\bigr) &= \bigl((w_1,0),0\bigr)\\
\bar f^{\#}\bigl((0,1),0\bigr) &= \bigl((0,w_2),0\bigr)\\
\bar f^{\#}\bigl((0,0),1\bigr) &= \bigl((0,0),1\bigr),
\end{aligned}
\tag{4.5}
\end{equation}

% [Tutor Note: 原文 p.139–140，式 (4.2)–(4.5)。核对锚点：$S_e = \mathbb{N}^2\oplus_{\mathbb N}\mathbb{N}$ 的两条
%  结构映射为对角映射 $\mathbb{N}\to\mathbb{N}^2$ 与乘 $e$ 的 $\mathbb{N}\to\mathbb{N}$，因此在 $S_e$ 中有
%  $((0,0),e)=((1,1),0)$；类似地在 $S_\ell$ 中 $((0,0),\ell)=((1,1),0)$。这正是把 (4.5) 的三条代入后推出
%  $w_1=w_2=\ell/e$ 的唯一依据（见下页）。$(4.3)$ 中左右两条竖箭头分别是 $1\mapsto\bar\rho_C$ 与 $1\mapsto\rho$，
%  它们是鬼层截面在茎上的取值。]

% [原文 p.141]
最后一条由 (4.3) 的交换性得到。但由于在 $S_e$ 中 $((0,0),e) = ((1,1),0)$，而在 $S_\ell$ 中
$((0,0),\ell) = ((1,1),0)$，我们在 $S_e$ 中有
\[
\bigl((\ell,\ell),0\bigr) = \bigl((0,0),\ell e\bigr)
= \bar f^{\#}\bigl((0,0),\ell e\bigr) = \bar f^{\#}\bigl((e,e),0\bigr)
= \bigl((ew_1,ew_2),0\bigr).
\]
由此断定 $w_1 = w_2 = \ell/e$。

这就给出命题的其余部分。
\end{proof}

理解了环面横截对数曲线的这些限制之后，我们现在可以定义与之相伴的热带曲线。

\begin{definition}
设 $f : C^{\dagger} \to X_0^{\dagger}$ 是 $\operatorname{Spec}k^{\dagger}$ 上的环面横截对数曲线。若 $C$ 的
两个不可约分支交于一个不映入 $\operatorname{Sing}(X_0)$ 的节点，则称它们是\emph{不可区分的}。

我们在 $C$ 的不可约分支之集上定义一个等价关系：称两个不可约分支 $C_1$ 与 $C_n$ \emph{等价}，如果存在一列
不可约分支 $C_1,C_2,\dots,C_n$，使得对 $1\le i\le n-1$，$C_i$ 与 $C_{i+1}$ 不可区分。

设 $\widetilde{\Gamma}_f$ 是（带若干非紧边、且可能带若干二价顶点的）带权图，满足：
\begin{enumerate}[label=(\arabic*),leftmargin=*]
\item $\widetilde{\Gamma}_f$ 的顶点集是 $C$ 的不可约分支模上述等价关系。若 $C'$ 是 $C$ 的不可约分支，我们用
$V_{C'}$ 记对应于 $C'$ 的等价类的顶点。
\item $\widetilde{\Gamma}_f$ 的有界边与映入 $\operatorname{Sing}(X_0)$ 的 $C$ 的节点一一对应。若 $y$ 是这样
一个节点，用 $E_y$ 记相应的边。若 $y\in C_1\cap C_2$，其中 $C_1,C_2$ 是 $C$ 的不同分支（由命题 4.9 这样的
分支存在），则 $E_y$ 的端点是 $V_{C_1}$ 与 $V_{C_2}$。若 $f(x)\in D_\omega$，其中 $\omega\in\cP$ 是仿射长度
为 $\ell$ 的边，且 $\overline{\cM}_{C,x} = S_e$，则边 $E_x$ 的权 $w(E_x)$ 为 $\ell/e$。由命题 4.9 (4)，这是
正整数，且等于该命题 (2) 中所定义的两个整数 $w_1,w_2$ 中的任一个。
\item $\widetilde{\Gamma}_f$ 的无界边与集合 $f^{-1}(\partial X_0)$ 一一对应，其中 $\partial X_0$ 由 (4.1)
定义。设 $p\in f^{-1}(\partial X_0)$ 对应于无界边 $E_p$。若 $p\in C'$，其中 $C'$ 是 $C$ 的不可约分支，则
$E_p$ 的端点是 $V_{C'}$，而权 $w(E_p)$ 是 $C'$ 在 $f(p)$ 处与 $\partial X_0$ 的相交重数。
\end{enumerate}
给定这个带权图，我们如下定义热带曲线
\[
h : \widetilde{\Gamma}_f \to M_\R
\]
(1) 若 $f(C')\subseteq D_v$，则 $h(V_{C'}) = v\in\cP$；（2）对 $y\in C_1\cap C_2$，$h(E_y)$ 是连接
$h(V_{C_1})$ 与 $h(V_{C_2})$ 的线段；（3）$h(E_p)$ 是 $\cP$ 中对应于 $\partial X_0$ 中含 $f(p)$ 的一维环面层
的边 $\omega$。
\end{definition}

注意 $\widetilde{\Gamma}_f$ 可能含二价顶点，而我们在热带曲线的定义（定义 1.11）中不允许这种顶点。这些二价顶点
来自那些只含两个在 $f$ 下映入 $\operatorname{Sing}(X_0)$ 的节点的不可约分支（或等价类）。要得到真正的热带
曲线，可以去掉这些二价顶点，把由二价顶点连接起来的一串边替换成

% [Tutor Note: 原文 p.141，定义 4.10 第 (2) 条。原文在这里写“$\overline{\cM}_{C,x}=S_e$ 时边 $E_x$ 的权
%  $w(E_x)$ 为 $\ell/e$”，但上下文中的节点记作 $y$、相应的边记作 $E_y$，紧接着又写“$E_x$ 的权”。这是原文把
%  $E_y$ 误写成 $E_x$（此处 $x$ 指的正是那个节点，故两种写法指向同一条边），译文照录不改；正确形式应为
%  $w(E_y)=\ell/e$。后面 Proposition 4.11 与 (2) 用的是 $E_y$。]

% [原文 p.142]
一条边。这给出新图 $\Gamma_f$，它与 $\widetilde{\Gamma}_f$ 有相同的同胚型。于是我们可以把 $h$ 看作给出映射
\[
h : \Gamma_f \to M_\R.
\]

有两件需要注意的事。第一，给定由二价顶点连接起来的一串边，需要验证这些边的权都相同，从而 $\Gamma_f$ 中对应于
这一串边的那条边有良定的权。第二，需要验证整条链在 $h$ 下的像是一条仿射线段，从而 $h$ 把 $\Gamma_f$ 中的新边
映成一条仿射线段。这两件事都可由 $\widetilde{\Gamma}_f$ 的平衡条件推出，见下面的命题。

\begin{proposition}
若 $f : C^{\dagger} \to X_0^{\dagger}$ 是环面横截对数曲线，则 $h : \Gamma_f \to M_\R$ 是参数化热带曲线。
\end{proposition}

\begin{proof}
我们先证明平衡条件（定义 1.11 的 (2)）对 $h : \widetilde{\Gamma}_f \to M_\R$ 成立。一旦知道这一点，那么
$\widetilde{\Gamma}_f$ 的二价顶点所相邻的任意两条边就有相同的权，并映入 $M_\R$ 中的同一条仿射线。这就说明
相应的 $h : \Gamma_f \to M_\R$ 现在是参数化热带曲线。事实上，如上面所讨论的，$\Gamma_f$ 的每条边现在都有
良定的权，并映入一条仿射线段。$\Gamma_f$ 中的任何顶点本来已是 $\widetilde{\Gamma}_f$ 中的顶点，因此平衡条件
对 $h : \Gamma_f \to M_\R$ 仍成立。

为证明 $\widetilde{\Gamma}_f$ 的平衡性，设 $V$ 是一个顶点，$C'\subseteq C$ 是对应于该顶点的等价类中诸不可约
分支之并。必有 $h(C')\subseteq D_v$ 对某个顶点 $v\in\cP$ 成立。设 $\Sigma_v$ 是 $M_\R$ 中定义 $D_v$ 的扇，
$m_1,\dots,m_r\in M$ 是 $\Sigma_v$ 的各射线的本原生成元，它们对应于 $D_v$ 的环面除子 $D_1,\dots,D_r$。这些
除子在 $D_v$ 的零维环面层之外是 Cartier 的，因此 $(f|_{C'})^*(D_i)$ 是 $C'$ 上的 Cartier 除子。令
\[
w_i := \deg(f|_{C'})^*(D_i).
\]
显然 $w_i$ 是 $\widetilde{\Gamma}_f$ 中与 $V$ 相邻、映入 $V+\R_{\ge0}m_i$ 的各边权之和（这是由权的定义）。
于是我们需要证明 $\sum_{i=1}^r w_im_i = 0$。为证这一点，我们对任意 $n\in N$ 证明
\[
\sum_{i=1}^r w_i\langle n,m_i\rangle = 0.
\]
注意 $\langle n,m_i\rangle$ 是 $z^n$ 在 $D_i$ 上的零点阶，故 $\sum w_i\langle n,m_i\rangle$ 是
$(f|_{C'})^*(z^n)$ 在 $C'$ 上的零点与极点的个数。当然它为零。
\end{proof}

% [Tutor Note: 原文 p.142，命题 4.11 的证明。核对锚点：$w_i$ 的两种读法必须一致——几何上它是 $\Gamma_f$ 中
%  映入方向 $m_i$ 的边权之和，代数上它是除子 $(f|_{C'})^*(D_i)$ 的次数。平衡条件 $\sum w_im_i=0$ 等价于
%  对偶形式 $\sum w_i\langle n,m_i\rangle=0$（$n$ 取遍 $N$），而后者就是“$z^n$ 在紧曲线 $C'$ 上零点数=极点数”。
%  这一步是整节的枢纽：它把对数曲线的组合数据与热带曲线的平衡条件接上。]

最后，我们考虑上述情形的一个装饰。回忆我们已固定点 $P_1,\dots,P_s \in M$ 与
$Q_1,\dots,Q_s \in \mathbb{G}(\widetilde{M})$，从而给出 $\pi$ 的截面 $\sigma_i$。令 $q_i = \sigma_i(0)$。

\begin{definition}
用 $(C^{\dagger},x_1,\dots,x_s)$ 记 $\operatorname{Spec}k^{\dagger}$ 上的一条对数曲线连同点
$x_1,\dots,x_s \in C$ 的选取，这些点都是光滑点（不是对数标记点）。则\emph{经过
$q_1,\dots,q_s$ 的环面横截带标记对数曲线}是指 $\operatorname{Spec}k^{\dagger}$ 上的对数态射
$f : (C^{\dagger},x_1,\dots,x_s) \to X_0^{\dagger}$，满足 $f$ 是 $X_0^{\dagger}$ 中的环面横截对数曲线，且
$f(x_i) = q_i$。
\end{definition}

给定这样的映射，我们如定义 4.10 那样联系于 $f$ 一个图 $\widetilde{\Gamma}_f$，但对每个 $x_i$ 额外加一条
无界边 $E_{x_i}$，且当

% [原文 p.143]
$x_i\in C'\subseteq C$ 时 $E_{x_i}$ 的端点为 $V_{C'}$。令 $w(E_{x_i}) = 0$，并完全像前面那样定义
$h : \widetilde{\Gamma}_f \to M_\R$，其中 $h|_{E_{x_i}}$ 为常值。

像前面一样，我们可以去掉剩余的二价顶点（在这个带标记的情形中可能更少，因为在无标记情形中本应是二价顶点的
地方，现在可能已经获得了一条相邻的、无界的标记边）。这给出
\[
h : \Gamma_f \to M_\R.
\]
我们还用
\[
h : \widehat{\Gamma}_f \to M_\R
\]
记由定义 4.10 联系于无标记映射 $f : C^{\dagger} \to X_0^{\dagger}$ 的热带曲线。它由 $\Gamma_f$ 去掉标记边
$E_{x_i}$ 并去掉由此产生的额外二价顶点得到。

最后，我们观察到 $h : \Gamma_f \to M_\R$ 现在是满足 $h(x_i) = P_i$ 的带标记热带曲线：

\begin{proposition}
(1) 记住 $P_i$ 是 $\cP$ 的顶点，我们有 $q_i \in \operatorname{Int}(D_{P_i})$。

(2) $h : \Gamma_f \to M_\R$ 是带标记参数化热带曲线，且 $h(E_{x_i}) = P_i$。
\end{proposition}

\begin{proof}
(1) 考虑由锥 $C(P_i)\in\Sigma_{\cP}$ 确定的仿射开子集 $X_{P_i}\subseteq X$。设
$n_1,n_2\in\widetilde{N}$ 是 $(P_i,1)^{\perp}$ 的一组基，并令 $n_3 = (0,0,1)\in\widetilde{N}$。则
$\pm n_1$、$\pm n_2$ 与 $n_3$ 生成 $C(P_i)^\vee\cap\widetilde{N}$，故
\[
X_{P_i} \cong \operatorname{Spec}k[x_1^{\pm1},x_2^{\pm1},x_3],
\]
其中 $x_i = z^{n_i}$。此外 $X_{P_i}\cap D_{P_i}$ 由 $x_3 = 0$ 给出。

现在设 $Q_i$ 在这些坐标下由 $Q_i = (a_1,a_2,a_3)$ 给出，$a_i\in k^{\times}$，$i=1,2,3$。则对
$z\in\mathbb{G}(L_i) = k^{\times}$，有 $z\cdot Q_i = (a_1,a_2,za_3)$。显然
$(a_1,a_2,0)\in\overline{\mathbb{G}(L_i)\cdot Q_i}\cap X_{P_i}$，故
$q_i = (a_1,a_2,0)\in D_{P_i}$。

(2) 若 $x_i\in C'$，则 $h(E_{x_i}) = h(V_{C'}) = v$ 对某个 $v\in\cP$ 成立，而此时当然有
$f(C')\subseteq D_v$。由于 $f(x_i) = q_i \in \operatorname{Int}(D_{P_i})$，必须有 $v = P_i$。
\end{proof}

\subsection{热带世界 $\to$ 对数世界}
我们继续使用数据 $P_1,\dots,P_s \in M$、好的格多面体分解 $\cP$ 以及一般点
$Q_1,\dots,Q_s \in \mathbb{G}(\widetilde{M})$。这给出点 $q_1,\dots,q_s \in X_0$，其中
$q_i = \overline{\mathbb{G}(L_i)\cdot Q_i}\cap X_0$。

现在设给定带标记热带曲线
\[
h : (\Gamma,x_1,\dots,x_s) \to M_\R
\]
满足 $h(x_i) = P_i$。

我们想找出所有经过 $q_1,\dots,q_s$、且相伴热带曲线为 $h$ 的对数曲线
$f : (C^{\dagger},x_1,\dots,x_s) \to X_0^{\dagger}$。一般来说，这样的曲线 $h$ 可以相当复杂。曲线的每个顶点都
可以用各种亏格的曲线之并来代表，而这些曲线可以成族地移动。然而最简单的情形是 $h$ 为定义 1.19 意义下的单曲线。
这时我们可以精确得多。

于是现在假定 $h$ 是亏格零的单曲线。为确保这样的曲线只有有限多条，我们照例假定 $s = |\Delta|-1$，其中

% [原文 p.144]
$\Delta$ 是 $h$ 的次数。由于 $\cP$ 包含所有经过 $P_1,\dots,P_s$ 的有理热带曲线，$h$ 的像含于 $\cP$ 的
一维骨架中。

正如上一节所指出的，联系于对数曲线的热带曲线的初始构造给出的图可能带二价顶点。通过细分 $\Gamma$ 的边得到图
$\widetilde{\Gamma}$ 以及 $h : \widetilde{\Gamma} \to M_\R$，我们可以假定对点
$y\in\widetilde{\Gamma}$，$h(y)$ 是 $\cP$ 的顶点当且仅当 $y$ 是 $\widetilde{\Gamma}$ 的顶点，或 $y$ 含于一条
标记无界边中。当然 $\widetilde{\Gamma}$ 现在可能带二价顶点。

本节的主要定理是：

\begin{theorem}
在上述情形下，设对 $\widetilde{\Gamma}$ 的每条边 $E$，$h(E)$ 的仿射长度都被权 $w(E)$ 整除。则满足
$f(x_i) = q_i$、且相伴热带曲线为 $h$ 的环面横截带标记对数曲线
\[
f : (C^{\dagger},x_1,\dots,x_s) \to X_0^{\dagger}
\]
的条数为 $\operatorname{Mult}(h)$。
\end{theorem}

\begin{remark}
关于每条边仿射长度可除性的假设，由联系于对数曲线的热带曲线的构造可知是必要的。另一方面，把格 $M$ 再作一次
重标度，总可以实现这种可除性。稍后我们会把这一步重标度说得更精确。
\end{remark}

在着手证明之前，我们先考察 $C$ 的不可约分支的性质，并注意 $\Gamma$ 的单性蕴含 $\widetilde{\Gamma}$ 的所有顶点
要么是二价的、要么是三价的。

\begin{definition}
设 $Y$ 是完整环面曲面。$Y$ 上的\emph{直线}是指一个非常值、环面横截的映射 $\phi : \PP^1 \to Y$，使得
$\#\phi^{-1}(\partial Y)\le3$，且对 $Y$ 的任意环面除子 $D$ 有 $\#\phi^{-1}(D)\le1$。
\end{definition}

给定一条直线，固定如下记号。设 $u_1,\dots,u_p\in M$ 是定义 $Y$ 的扇中、对应于那些满足
$\phi^{-1}(D_i)\ne\emptyset$（$1\le i\le p$）的环面除子 $D_1,\dots,D_p\subseteq Y$ 的各射线的本原生成元。
此外设 $w_i$ 是 $\phi^*(z_i)$ 的零点阶，其中 $z_i$ 是 $D_i$ 的局部方程。这给出 $\phi$ 与 $D_i$ 的相切阶。
于是如命题 4.11 的证明中那样，平衡条件成立：
\[
\sum_i w_iu_i = 0.
\]
由于 $p\le3$，这特别蕴含 $p = 2$ 或 $3$。我们分别称这两种情形为\emph{二价情形}与\emph{三价情形}。

% [Tutor Note: 原文 p.144，定义 4.16 之后的论证。“$p=2$ 或 $3$”可直接核对：若 $p=1$，平衡条件给出
%  $w_1u_1=0$，而 $u_1\ne0$ 是本原向量，故 $w_1=0$；但 $w_1$ 是 $\phi$ 与 $D_1$ 的相切阶，$\phi$ 非常值且
%  环面横截，必有 $w_1>0$，矛盾。所以 $p\ge2$，再由 $p\le3$ 得到结论。另注意“直线”在这里指的是映射
%  $\phi:\PP^1\to Y$ 而非 $Y$ 的子簇，这与热带几何中“热带直线”是参数化曲线而非子集是同一精神。]

称带有相联数据 $\mathbf{u} = (u_1,u_2,\dots)$、$\mathbf{w} = (w_1,w_2,\dots)$ 的直线为
$(\mathbf{u},\mathbf{w})$ 型的，并用 $\mathcal{L}_{(\mathbf{u},\mathbf{w})}$ 记环面曲面 $Y$ 中这种类型的
所有直线之集（模去定义域 $\PP^1$ 的自同构）。我们现在来分类直线，先处理二价情形：

\begin{lemma}
设 $(\mathbf{u},\mathbf{w}) = ((u_1,u_2),(w_1,w_2))$，并设 $E$ 是由 $u_1$ 生成的 $M$ 的一维子格（由平衡
条件，等价的也可用 $u_2$ 生成）。存在由格映射 $M \to M/E$ 诱导的自然映射 $g : Y \to \PP^1$。则每个
$\phi\in\mathcal{L}_{(\mathbf{u},\mathbf{w})}$ 的像为 $g^{-1}(p)\subseteq Y$，其中 $p\in\PP^1\setminus\partial\PP^1$；
并且映射 $\phi : \PP^1 \to g^{-1}(p)$ 是 $g^{-1}(p)$ 的 $w_1$ 重分歧覆盖，恰在两点
$g^{-1}(p)\cap\partial Y$ 处完全分歧。
\end{lemma}

% [原文 p.145]
\begin{proof}
由于 $\rho_i = \R_{\ge0}u_i$ 是定义 $Y$ 的扇 $\Sigma_Y$ 中的锥，且 $u_2 = -u_1$，$\Sigma_Y$ 中每个锥都含于
包含 $\R u_i$ 的两个半平面之一，从而映入 $(M/E)\otimes_\Z\R$ 中的一条射线。于是投影 $M\to M/E$ 定义了从
$\Sigma_Y$ 到 $(M/E)\otimes_\Z\R$ 中定义 $\PP^1$ 的扇的扇映射。于是我们可以把 $Y$ 的环面边界写成
$\partial Y = D_1\cup D_2\cup g^{-1}(0)\cup g^{-1}(\infty)$。由于 $\phi(\PP^1)$ 与 $g^{-1}(0)$ 及
$g^{-1}(\infty)$ 不相交，我们看到 $\phi(\PP^1)$ 必是 $g$ 的一条纤维，但不是 $g^{-1}(0)$ 或 $g^{-1}(\infty)$。
$g$ 的所有纤维都同构于 $\PP^1$，$g^{-1}(0)$ 与 $g^{-1}(\infty)$ 除外。要使 $\phi^{-1}(\partial Y)$ 只由两点
组成，$\phi$ 必须在 $\partial Y$ 处完全分歧，于是得到结论。
\end{proof}

\begin{corollary}
在引理 4.17 的情形下，$\mathcal{L}_{(\mathbf{u},\mathbf{w})} \cong \mathbb{G}(M/E)$。
\end{corollary}

\begin{proof}
这不过是观察到引理中所述的覆盖是唯一的，因此 $\phi\in\mathcal{L}_{(\mathbf{u},\mathbf{w})}$ 完全由一点
$p\in\mathbb{G}(M/E) = \PP^1\setminus\partial\PP^1$ 确定。
\end{proof}

接下来考虑三价情形，于是固定
$(\mathbf{u},\mathbf{w}) = ((u_1,u_2,u_3),(w_1,w_2,w_3))$。为方便起见，可以假定定义 $Y$ 的扇
$\Sigma_Y$ 只有三条射线，由 $\rho_i = \R_{\ge0}u_i$ 定义。事实上，总存在到这样的曲面的环面收缩
$\alpha : Y \to Y'$；且由环面横截性，$Y$ 与 $Y'$ 上 $(\mathbf{u},\mathbf{w})$ 型直线之集相同。

\begin{lemma}
在上述情形下，设 $\Sigma_Y$ 是由 $u_1,u_2,u_3$ 生成的、恰有三条射线的完整扇。考虑映射
\[
\Z^2 \to M
\]
由
\[
(a,b) \mapsto aw_1u_1+bw_2u_2
\]
给出。它诱导从例 1.14 所给出的 $\PP^2$ 的标准扇 $\Sigma_{\PP^2}$（在 $\R^2$ 中）到 $\Sigma_Y$ 的扇映射，从而
诱导映射
\[
f_{(\mathbf{u},\mathbf{w})} : \PP^2 \to Y.
\]
$(\mathbf{u},\mathbf{w})$ 型的任何直线 $\phi : \PP^1 \to Y$ 都同构于一个环面横截线性嵌入
$\PP^1 \to \PP^2$ 与 $f_{(\mathbf{u},\mathbf{w})}$ 的复合。
\end{lemma}

\begin{proof}
映射 $\Z^2 \to M$ 显然把第一象限送到 $u_1$ 与 $u_2$ 张成的锥。注意由平衡条件该映射还把 $(-1,-1)$ 送到
$-w_1u_1-w_2u_2 = w_3u_3$，由此可知 $\Sigma_{\PP^2}$ 中由 $(1,0)$ 与 $(-1,-1)$ 生成的锥被送到 $\Sigma_Y$ 中
由 $u_1$ 与 $u_3$ 生成的锥，而 $\Sigma_{\PP^2}$ 中由 $(0,1)$ 与 $(-1,-1)$ 生成的锥被送到 $\Sigma_Y$ 中由 $u_2$
与 $u_3$ 生成的锥。这给出扇映射，从而给出 $f_{(\mathbf{u},\mathbf{w})} : \PP^2 \to Y$。

现在设 $C$ 是 $\PP^1\times_Y\PP^2$ 的某个不可约分支的正规化，其中映射 $\PP^1\to Y$ 为 $\phi$、映射
$\PP^2\to Y$ 为 $f_{(\mathbf{u},\mathbf{w})}$。我们有投影 $\tilde\phi : C \to \PP^2$。只需说明：第一，$\tilde\phi$
是 $\PP^2$ 中一条直线的嵌入；第二，复合 $C \to \PP^1\times_Y\PP^2 \to \PP^1$ 是同构。

首先注意 $f_{(\mathbf{u},\mathbf{w})}$ 的次数等于大环面轨道映射 $\mathbb{G}(\Z^2)\to\mathbb{G}(M)$ 的次数，
而这正是 $\Z^2$ 在 $M$ 中的像的指数，即由 $w_1u_1$ 与 $w_2u_2$ 生成的 $M$ 的子格的指数 $\delta$。

我们还计算 $f_{(\mathbf{u},\mathbf{w})}$ 沿除子 $D_i$ 的分歧次数。不失一般性，考虑除子 $D_1$，并设在 $M$ 的
某组基下 $u_1 = (1,0)$、

% [原文 p.146]
$u_2 = (a,b)$。设 $\rho_i$ 是由 $u_i$ 生成的锥。则 $\rho_1^\vee$ 由 $(0,\pm1)$ 与 $(1,0)$ 生成，而
$f_{(\mathbf{u},\mathbf{w})}$ 在由 $\R_{\ge0}(1,0)$ 与 $\rho_1$ 分别确定的 $\PP^2$ 与 $Y$ 的仿射卡上，由环映射
\[
k[x,y^{\pm1}] \to k[x,y^{\pm1}]
\]
给出，定义为
\[
x \mapsto x^{w_1}y^{w_2a},\qquad y \mapsto y^{w_2b}.
\]
由于 $y$ 是单位，而除子 $D_1$ 由 $x=0$ 给出，我们看到 $f_{(\mathbf{u},\mathbf{w})}$ 在 $D_1$ 上有 $w_2b$ 个
分支，且每个分支在 $D_1$ 上的分歧阶为 $w_1$。更一般地，由于 $f_{(\mathbf{u},\mathbf{w})}$ 的次数为 $\delta$
（在此情形为 $bw_1w_2$），我们看到 $f_{(\mathbf{u},\mathbf{w})}$ 在 $D_i$ 上有 $\delta/w_i$ 个分支，每个分支
的分歧阶为 $w_i$。注意这个分歧阶与 $\phi$ 在该唯一交点处与 $D_i$ 的相交指数一致。

现在考虑图

\begin{center}
\includegraphics[width=0.60\textwidth]{f4_diagram2.png}
\end{center}

% [Tutor Note: 原文 p.146，引理 4.19 证明中的交换图（原书无编号）。此图箭头交叉较多（$C$ 同时经
%  $\tilde\phi$ 映到 $\PP^2$、经 $\psi$ 映到 $\PP^1\times_Y\PP^2$，另有 $\pi_1,\pi_2,\varphi,f_{(\mathbf u,\mathbf w)}$），
%  用 array 排会严重失真，故按 STYLE_GUIDE §5“示意图优先按原图裁剪”处理，字体差异由裁剪保留原貌。
%  核对锚点：图中方形 $\PP^1\times_Y\PP^2 \to \PP^2 \to Y$ 与 $\PP^1\times_Y\PP^2 \to \PP^1 \to Y$ 两个方向都
%  复合到 $Y$；$\psi$ 是 $C$ 到纤维积的态射，$\tilde\phi$ 是它在 $\PP^2$ 上的投影。]

局部地，比如在 $D_1$ 附近，在上述坐标下有 $\phi^*(x) = \text{单位}\cdot t^{w_1}$，其中 $t$ 是 $\PP^1$ 在
$\phi^{-1}(D_1)$ 某邻域内的局部参数，且在 $\phi^{-1}(D_1)$ 处为零。另一方面，
$f^*_{(\mathbf{u},\mathbf{w})}(x) = x^{w_1}\cdot y^{w_2a} = x^{w_1}\cdot\text{单位}$，故
$t^{w_1} = x^{w_1}\cdot\text{单位}$ 是 $\PP^1\times_Y\PP^2$ 在 $\PP^1\times\PP^2$ 中、在点 $p$ 某邻域内的
理想的一个局部方程，其中 $p\in\PP^1\times_Y\PP^2$ 满足 $\phi(\pi_1(p))\in D_1$。由此我们看到有 $w_1$ 个
$\PP^1\times_Y\PP^2$ 的分支经过 $p$，每个分支经 $\pi_1$ 局部同构地映到 $\PP^1$。正规化之后，这 $w_1$ 个分支
被分开，所以若 $p'\in C$ 满足 $\psi(p') = p$，则 $\psi$ 在 $p'$ 的某邻域内非分歧。我们还看到，若 $t'$ 是 $C$
在 $p'$ 邻域内的局部参数，则 $\tilde\phi^*(x) = \text{单位}\cdot t'$。于是 $\tilde\phi$ 与
$f^{-1}_{(\mathbf{u},\mathbf{w})}(D_1)$ 横截。同样的分析在 $D_2$ 与 $D_3$ 处也成立，于是我们看到 $\psi$ 非分歧，
从而由于 $\PP^1$ 单连通，$\psi$ 是同构。此外 $\tilde\phi$ 与 $\PP^2$ 的每个环面除子横截，因此是一条直线的嵌入。
\end{proof}

\begin{corollary}
在引理 4.19 的情形下，$\mathbb{G}(M)$ 在 $\mathcal{L}_{(\mathbf{u},\mathbf{w})}$ 上单传递地作用。特别地，
$\mathcal{L}_{(\mathbf{u},\mathbf{w})}$（非典范地）同构于 $\mathbb{G}(M)$。
\end{corollary}

\begin{proof}
任何直线 $\phi : \PP^1 \to Y$ 至多提升为 $\PP^2$ 中 $\delta = \deg f_{(\mathbf{u},\mathbf{w})}$ 条不同的直线，
而 $\delta$ 是同态 $\mathbb{G}(\Z^2)\to\mathbb{G}(M)$ 的核的阶。现在 $\mathbb{G}(\Z^2)$ 显然在环面横截线性
嵌入 $\tilde\phi : \PP^1 \to \PP^2$ 之集上单传递地作用，于是
$\ker(\mathbb{G}(\Z^2)\to\mathbb{G}(M))$ 在 $\phi$ 到 $\PP^2$ 的各个提升 $\tilde\phi$ 之集上传递地作用。
这说明 $\mathbb{G}(M)$ 在 $\mathcal{L}_{(\mathbf{u},\mathbf{w})}$ 上的作用是单传递的。
\end{proof}

现在我们将把这一切汇总起来，构造所需对数曲线的底层概形态射 $f : C \to X_0$。我们将称这样的曲线为
\emph{前对数曲线}：

% [Tutor Note: 原文 p.144–146，§4.3 前半（定理 4.14、定义 4.16、引理 4.17/4.19、推论 4.18/4.20）。
%  核对锚点一：定理 4.14 要求“每条边的仿射长度被权整除”，这正是命题 4.9 (4) 的 $ew_i=\ell$ 中 $e\in\Z_{\ge1}$
%  的等价说法；故 Remark 4.15 说重标度 $M$ 总能实现它：把 $M$ 换成超格 $\frac1p M$ 时本原向量变短，
%  同一条线段的仿射长度按 $p$ 倍放大，取 $p$ 为各权 $w(E)$ 的公倍数即可使 $w(E)\mid\ell(E)$。
%  核对锚点二：引理 4.19 的次数计算 $\deg f_{(\mathbf{u},\mathbf{w})} = \delta = \operatorname{index}(w_1u_1,w_2u_2)$
%  与“$D_i$ 上 $\delta/w_i$ 个分支、每支分歧阶 $w_i$”相容，因为 $\delta = w_i\cdot(\delta/w_i)$。
%  本块在此处结束（“…pre-log curves:”）；Definition 4.21 起属 c04b。]
